\scanpage{109}
\mspage{100}
% archive sheet number: 109

et $\Psi$ soient les comp. connexes de $\unsure{\Phi|\varphi}\cup\Psi$, i.e. $\Phi\cup\Psi$ disconnexe, il f. et s. que $\Phi\overset{\unsure{!|'|1}}{<}\Psi$ ou $\Psi\overset{!}{<}\Phi$, i.e.
\[
\{\Phi,\Psi\}\in\operatorname{Drap}_{<^!}(\operatorname{Préfig}_{\mathrm{conn}}(\unsure{\mathscr{M}|\mathcal{M}})).
\]

En résumé, on trouve donc la petite proposition récapitulative

\uline{Théorème} 1) Pour tout $\unsure{\mathbb{D}|D}\in\operatorname{Drap}^{*}(L)$, et toute partie $\alpha$ de $\unsure{\partial|\delta}\mathbb{D}$ (cf\unsure{...}), soit $\Phi_{\mathbb{D},\alpha}=F_{\mathbb{D}}\setminus\alpha$, où $F_{\mathbb{D}}$ est la figure formée des strates \unsure{interstitielles} du drapeau $D$.
\margin{$\partial D=\left\{\begin{aligned}&\varnothing\text{ si }\operatorname{card}D=1\\&\{\operatorname{or}D,\operatorname{ex}D\}\text{ si}\\&\quad\operatorname{card}D\geq2\end{aligned}\right.$}

On trouve ainsi une bijection
\[
(\mathbb{D},\alpha)\unsure{\longmapsto|\longrightarrow}\Phi_{\mathbb{D},\alpha}
\]
de l'ens des paires $D,\alpha$ sur l'ens des préfigures connexes.

2) Pour que $\Phi$ soit une figure i.e. fermée, il f. et s. que $\alpha=\varnothing$. Pour que $\Phi$ soit ouverte, \deleted{\unsure{et on a}} il f. et s., (en admettant Hyp) \deleted{\unsure{que et}} et en supposant que $L$ n'a pas de plus petit ni de plus grand élément, que \deleted{$\alpha$} $\operatorname{card}\mathbb{D}\geq2$ et $\alpha=\partial\mathbb{D}$. En \deleted{\unsure{disposant par}} admettant Hyp. et sans \unsure{préjudice|préjuger} s'il existe un plus petit ou
